\chapter{讨\quad 论}\label{chap:discussion}

\section{主要发现}

\subsection{喷嘴几何参数对打印性能的影响机制}

本研究系统考察了四个喷嘴与整形板几何参数（$\theta$、$L$、$TL$、$\alpha$）对基底接触压力及相关打印性能指标的影响，结果揭示了各参数作用机制的本质差异，为打印喷嘴的工程设计提供了清晰的物理背景。

\textbf{（1）喷嘴收缩角$\theta$的双重作用机制}

收缩角$\theta$通过两条相互独立的物理路径影响打印性能。其一是压力路径：较大的收缩角使流道截面积沿轴向更急剧地减小，浆体在收缩段内受到更强的径向压缩，出口截面动压更高，经前挡板阻挡后转化为更大的基底接触压力，为下游整形板区域提供更高的压力输入。其二是剪切路径：更大的收缩角使流线沿收缩壁面的曲率变化趋于平滑，抑制了收缩段转角处的流动分离，消除了局部高剪切速率集中区，有效降低壁面剪切应力。这两条路径的方向一致，赋予了收缩角参数独特的优化特征，在不引入任何指标劣化的前提下同时改善接触压力和剪切应力，具有重要的工程价值。

\textbf{（2）喷嘴竖直高度$L$的有限调控作用}

相比$\theta$和整形板参数，喷嘴竖直高度$L$对基底接触压力的影响最为有限。从流动的角度理解，喷嘴竖直高度的主要功能是为流体提供充分发展的速度剖面，保证进入下游收缩段的流动状态稳定，而非主动产生动压或楔形约束效应。延长直管段仅增加管道摩擦损失，对出口截面的速度分布和动压幅值几乎无影响，因而对基底接触压力的调控能力极弱。整形板末端剪切应力 $WSS_{outlet}$ 随$L$变化的先减后升则反映了竖直高度通过影响速度剖面充分发展程度而间接作用于整形板区域的局部流动状态，很可能属于二阶效应。综合来看，$L$在本文研究范围内是一个比较弱的影响参数，在满足制造约束的前提下，其具体取值对打印性能的影响可后置考虑。

\textbf{（3）整形板结构的不可替代性}

无顶板对照工况（C2）的CFD结果说明：去除整形板顶板后，多项指标成为全部工况中最低值，甚至低于无整形板的理论自由挤出状态。这一结果表明，整形板顶板并非简单的成形模具，而是产生基底接触压力的核心功能构件。其作用机制在于：顶板封闭了流体向上方逃逸的路径，将喷嘴出口的高动压流体强制约束在楔形通道内，使动能得以有效转化为作用于基底面的静压。一旦顶板缺失，流体以近乎无约束的方式向四周扩散，原本用于产生接触压力的动能大量耗散，导致整体接触压力大幅下降，也会导致整体挤出性能下降。

\textbf{（4）整形板参数协同效应}

C1-3工况的 $P_{sub,AWA}$ 提升幅度显著超过各单因素效果的线性叠加，揭示了喷嘴参数与整形板参数之间存在正向协同增强机制。$\theta=60°$的大收缩角使喷嘴出口流体具有更高的动压和更均匀的速度剖面，而$\alpha=15°$与$TL=60$~mm所形成的强楔形约束空间对高动压流体的压力转化效率更高。在这个体系下，输入动压越高，楔形效应产生的压力增益越大，两种效应的协同效应导致总压力增益超过各自独立贡献之和。这一协同机制的存在意味着，在实际打印系统优化中，喷嘴参数（$\theta$）与整形板参数（$\alpha$、$TL$）更应作为一个整体系统加以联合优化，而非独立调节。

\subsection{打印性能指标之间的耦合与权衡关系}

\textbf{（1）压实性能与泵送能耗的高度共线性}

全部20个工况的 $P_{inlet}$ 与 $P_{sub,AWA}$ 之间的Pearson相关系数高达 $r=0.994$（$p<0.001$），两者近乎完全共线，线性拟合关系为 $P_{inlet} \approx 1.77 \times P_{sub,AWA}$。这一发现揭示了当前喷嘴-整形板系统的一个根本性约束：提升层间压实效果与降低泵送能耗是相互对立的，在本课题的研究系统架构下无法独立解耦。本研究所采用的几何参数（$\theta$、$L$、$TL$、$\alpha$）均通过改变流动阻力来调控接触压力，而增大接触压力必然要求克服更大的流动阻力，更大的流动阻力直接体现为更高的入口驱动压力。这一内在耦合的存在，在五目标Pareto分析中导致 $P_{inlet}$ 与 $P_{sub,AWA}$ 成为近乎等价的目标维度，也就导致了完全共线的拟合关系。在工程实践中，若想要提高基底接触面的压力大小，则必须需要这必须在系统设计阶段预先考量泵送系统的工作压力限制，找到一个系统平衡点。

\textbf{（2）压实性能与剪切应力的条件性权衡}

在一般情况下，提升接触压力往往伴随着更强的流场约束和更大的速度梯度，导致壁面剪切应力升高，但本研究下喷嘴出口倾斜角度$\theta$同时实现了接触压力提升和壁面剪切应力降低。另外，倒角处理则提供了另一条打破常规权衡的途径：在不牺牲接触压力的前提下，通过消除尖角处的流动分离，倒角可以显著降低壁面剪切应力。因此，压实与剪切之间的权衡关系并非不可调和，而是取决于具体的几何优化手段。宏观喷嘴几何和局部转角几何均可在特定条件下同时改善两个目标。

\textbf{（3）Pareto Front的几何形态与设计含义}

在以最大化基底接触压力和最小化壁面剪切应力构成的二维目标空间中，CFD工况点的分布呈现出清晰的Pareto Front结构。C3工况位于四目标Pareto Front的唯一最优点，而C1-3仅因壁面剪切应力劣于C3而被支配。这一结构表明，在可行设计空间内，C3代表了当前参数组合下压实性能与剪切损伤控制的最优平衡点，继续改变任何单一参数均无法在不损失其他目标的前提下进一步改善综合性能。

NSGA-II在代理模型上搜索到的理论最优点（$TL=70$，$\theta=60°$，$\alpha=14.06°$，$H=5$~mm）虽然能将基底接触压力提升超过一倍，但其泵送压力也同步提升了一倍，压力传递效率 $\eta_p$ 仅仅只有微弱的提升，清晰地说明了在层高约束的硬性边界附近，继续追求极限压实性能的边际效益急剧下降，工程代价大幅上升，C3才是综合考量下最具工程实用价值的设计方案。

\subsection{层高约束对混凝土打印性能的主导效应}

在本研究的所有参数化发现中，\cref{eq:H}揭示的层高解析关系：
\begin{equation}
  H = Z_0 - TL \cdot \sin(\alpha) \tag{2.1}
\end{equation}
是贯穿全部结果的核心约束机制，其重要性远超单一几何参数的影响。从物理层面上来看，$\alpha$和$TL$对基底接触压力的影响并非独立的几何效应，而是两个共享同一物理中间量（混凝土层高 $H$）的耦合路径：$\alpha$增大通过增大 $\sin(\alpha)$ 来减小 $H$，$TL$增大通过增大前乘系数来放大 $\sin(\alpha)$ 的效果，两者都通过缩小出口间隙来增强Reynolds楔形压力，本质机制相同，只是调控路径不同。

这个约束具有重要的简化意义：在参数优化中，$\alpha$和$TL$并不是两个真正独立的自由度，而是在层高约束的框架下共同决定出口间隙高度的两个相关参数。当 $H$ 被固定时，$\alpha$和$TL$之间存在确定的几何替代关系，增大一个参数可以在保持 $H$ 不变的前提下等效替代另一个参数对接触压力的贡献。也正因如此，在NSGA-II的优化搜索中，$\alpha$和$TL$均趋向其各自的上界，共同将层高压缩至约束下界，这实际上正是层高约束作为主导变量主动锁定了最优解的位置。

如果要在确定$\alpha$和$TL$的基础上进一步提升接触压力而不受层高约束限制，有效的系统级改进路径有：增大提高喷嘴离地高度或减小 $H_{min}$。通过改进浆体配比提高早龄期强度，使更薄的沉积层仍能保持稳定，可能是打印混凝土性能提高的另一方向。

\section{工程设计建议}

综合参数化研究与多目标优化结果可知，不同几何参数对打印性能具有明显不同的作用机制。其中，整形板倾斜角 $\alpha$ 对基底接触压力的调控最为敏感，但同时会显著压缩层高余量；整形板长度 $TL$ 能够以较为线性且可预测的方式提升压实能力，同时对层高影响相对温和；收缩角 $\theta$ 在提高接触压力的同时，还能够显著降低喷嘴壁面剪切应力，表现出较好的协同优化特性。相比之下，而喷嘴竖向高度 $L$ 对整体性能影响较弱，在本研究范围内可视为次要参数。此外，转角圆角处理虽然对接触压力改善有限，但能够明显降低局部壁面剪切应力与潜在磨损风险，是一种具有较高工程性价比的优化措施。

综合五目标Pareto Front分析、TOPSIS排序以及流场结构分析结果，推荐采用C3工况作为最优设计方案，即 $\theta=60°$、$TL=60$~mm、$L=20$~mm、$\alpha=15°$，并在关键转角区域引入3--5~mm的圆角过渡处理。这个方案在保持可接受层高余量的同时，实现了较高的基底接触压力、较低的局部剪切损伤以及较优的压力传递效率。虽然入口压力同步增加约，但整体仍处于工程可接受范围内，表现出较优的综合打印性能。

\section{方法局限性讨论}

\subsection{CFD模型假设局限}

本研究的CFD模型建立在若干理想化假设基础上，以提高计算效率并控制参数复杂度，因此对真实3D打印混凝土过程的描述仍存在一定局限性。本课题采用稳态单相流模型，并没有考虑打印混凝土的触变性、结构重建以及多层沉积过程中的时变行为，而实际浆体的早龄期强度会随时间持续演化，并对层间稳定性与粘结性能产生重要影响。此外，模型将浆体视为均质连续介质，忽略了润滑层形成、骨料迁移及离析等多相流效应，这可能导致对壁面剪切应力和系统压降的预测偏高。

混凝土浆体在挤出时同样伴随着水化放热反应，但是CFD模拟假设等温条件与刚性光滑壁面，不仅无法模拟混凝土浆体水化热造成的温度变化，也无法模拟壁面粗糙度与磨损对流场的影响。另外，喷嘴出口与基底面之间的高度$Z_0$也是影响混凝土挤出性能的一大影响因素，在本研究中保持固定，其对接触压力与层高分布的独立影响尚未展开系统研究。未来工作可进一步引入触变性本构模型与瞬态多层打印模拟，以更全面地描述实际混凝土打印过程中的复杂流变行为与结构演化机制。

\subsection{代理模型精度局限}

本研究使用17个可行CFD工况数据点构建了4维多项式RSM代理模型，用于NSGA-II多目标优化搜索。2阶多项式RSM在4个设计变量下展开后包含15个基函数项，17个数据点仅提供了2个自由度用于残差估计，模型处于近过拟合状态，LOO交叉验证的预测误差必然偏大，泛化能力有限。在后续完善中，可以试图补充CFD数据点到45--70个，以便代理模型能够正常收敛，输出更加精准的预测值。

\subsection{实验验证缺失}

本研究为纯数值模拟研究，所有结论均基于CFD计算结果，尚缺乏实验室尺度打印试验的物理验证，这会导致模型预测结果在绝对数值层面存在一定的不确定性。打印混凝土的流变参数对模拟结果具有较高敏感性，本文采用的Herschel--Bulkley参数来源于文献中的典型打印混凝土数据，与实际材料配合比之间可能存在差异，这种差异可能进一步影响接触压力、壁面剪切应力以及系统压降的预测精度。

另外，基底接触压力在实验条件下较难直接测量，目前尚缺乏成熟且统一的测试方法，这也限制了本文关键评价指标的实验验证。未来工作可结合实验室打印平台，对代表性喷嘴工况开展入口压力、流量以及沉积形态等方面的对比测试，从而进一步评估本文CFD模型的工程适用性与预测可靠性。
